\documentclass{article}
\usepackage[UTF8]{ctex}
\usepackage{amsmath}
\usepackage{graphicx}
\usepackage{listings}
\usepackage{xcolor}
\usepackage{hyperref}
\usepackage{geometry}

% Page setup
\geometry{a4paper, margin=1in}

% Code listing setup
\lstset{
    backgroundcolor=\color{gray!10},
    basicstyle=\ttfamily\small,
    numbers=left,
    stepnumber=1,
    showstringspaces=false,
    tabsize=4,
    breaklines=true,
    frame=single
}

% Title setup
\title{软件工程文档标题}
\author{作者姓名}
\date{\today}

\begin{document}

\maketitle

\tableofcontents
\newpage

\section{引言}
这里是文档的介绍部分。

\section{主要内容}

\subsection{子章节}
这是子章节的内容。

\section{数学公式示例}
\label{sec:math}

\subsection{行内公式}
圆的面积公式为 $A = \pi r^2$，其中 $r$ 是半径。

\subsection{行间公式}
均方误差（MSE）定义为：
\begin{equation}
    \text{MSE} = \frac{1}{n} \sum_{i=1}^{n} (y_i - \hat{y}_i)^2
    \label{eq:mse}
\end{equation}

\subsection{多行公式}
\begin{align}
    \nabla \cdot \mathbf{E} &= \frac{\rho}{\varepsilon_0} \quad \text{(高斯定律)} \\
    \nabla \cdot \mathbf{B} &= 0 \quad \text{(磁场无散度)} \\
    \nabla \times \mathbf{E} &= -\frac{\partial \mathbf{B}}{\partial t} \quad \text{(法拉第定律)}
    \label{eq:maxwell}
\end{align}

\subsection{矩阵表示}
协方差矩阵：
\[
\Sigma = \begin{pmatrix}
\sigma_1^2 & \rho_{12}\sigma_1\sigma_2 & \rho_{13}\sigma_1\sigma_3 \\
\rho_{12}\sigma_1\sigma_2 & \sigma_2^2 & \rho_{23}\sigma_2\sigma_3 \\
\rho_{13}\sigma_1\sigma_3 & \rho_{23}\sigma_2\sigma_3 & \sigma_3^2
\end{pmatrix}
\]

\section{代码示例}
\label{sec:code}

Python 代码示例：
\begin{lstlisting}[language=Python]
def calculate_mse(y_true, y_pred):
    """计算均方误差"""
    n = len(y_true)
    mse = sum((y - y_hat) ** 2 for y, y_hat in zip(y_true, y_pred)) / n
    return mse
\end{lstlisting}

JavaScript 代码示例：
\begin{lstlisting}[language=JavaScript]
function fibonacci(n) {
    if (n <= 1) return n;
    return fibonacci(n - 1) + fibonacci(n - 2);
}
\end{lstlisting}

\section{表格示例}
\label{sec:table}

\begin{table}[htbp]
    \centering
    \caption{复杂度分析表}
    \begin{tabular}{|c|c|c|}
        \hline
        \textbf{算法} & \textbf{时间复杂度} & \textbf{空间复杂度} \\
        \hline
        冒泡排序 & $O(n^2)$ & $O(1)$ \\
        \hline
        快速排序 & $O(n \log n)$ & $O(\log n)$ \\
        \hline
        归并排序 & $O(n \log n)$ & $O(n)$ \\
        \hline
    \end{tabular}
\end{table}

\section{图片示例}
\label{sec:figure}

\begin{figure}[htbp]
    \centering
    \includegraphics[width=0.8\textwidth]{example-image}
    \caption{系统架构图}
    \label{fig:architecture}
\end{figure}

\section{结论}
\label{sec:conclusion}

本文档展示了 LaTeX 文档的各种特性，包括：
\begin{itemize}
    \item 数学公式（行内和行间）
    \item 编号公式和交叉引用
    \item 代码高亮
    \item 表格和图片
    \item 中文支持
\end{itemize}

\section{参考文献}

\begin{thebibliography}{99}
    \bibitem{ref1} Author Name. \textit{Book Title}. Publisher, Year.
    \bibitem{ref2} Author Name. ``Article Title''. \textit{Journal Name}, Year.
\end{thebibliography}

\end{document}
